
\documentclass[12pt]{article}

\usepackage[spanish]{babel}
\usepackage[utf8]{inputenc}
\usepackage{geometry}
\usepackage{longtable}
\usepackage{array}
\usepackage{titlesec}
\usepackage{setspace}

\geometry{letterpaper, margin=2.5cm}
\setstretch{1.5}

\title{Propuesta de desarrollo y despliegue de un proyecto de Inteligencia Artificial\\
\large Reconocimiento automático de vouchers de ventas por datáfono para cuadre de caja}
\author{}
\date{\today}

\begin{document}

\maketitle

\section{Introducción}

En los puntos de venta de tiendas y comercios minoristas, el proceso de cuadre de caja suele requerir la validación manual de vouchers generados por datáfonos. Este procedimiento implica registrar valores, fechas, números de autorización, franquicias y referencias de transacción dentro de una plataforma web administrativa.

Actualmente, muchos establecimientos realizan este procedimiento manualmente, lo que genera retrasos en el cierre de caja, errores de digitación, pérdida de vouchers físicos, inconsistencias entre las ventas registradas y los comprobantes bancarios, además de sobrecarga operativa para cajeros y auxiliares administrativos.

El presente proyecto propone el desarrollo de un sistema basado en Inteligencia Artificial y Visión por Computador capaz de reconocer automáticamente la información contenida en imágenes de vouchers de ventas por datáfono y registrar dichos datos directamente en la plataforma web utilizada para el cuadre de caja.

\section{Requisitos del sistema}

\subsection{Necesidad a resolver}

Automatizar el registro de información contenida en vouchers de ventas por datáfono para disminuir errores operativos y reducir el tiempo de cuadre de caja en los puntos de venta.

\subsection{Objetivos generales}

Diseñar e implementar una solución de Inteligencia Artificial capaz de capturar imágenes de vouchers, extraer automáticamente información relevante mediante OCR e IA, validar los datos obtenidos y registrar automáticamente la información en la plataforma web de cuadre de caja.

\subsection{Objetivos específicos}

\begin{itemize}
    \item Reducir al menos un 80\% del tiempo invertido en digitación manual.
    \item Disminuir errores humanos en el registro de vouchers.
    \item Permitir validaciones automáticas de valores y fechas.
    \item Centralizar digitalmente los comprobantes.
    \item Facilitar auditorías y trazabilidad financiera.
\end{itemize}

\section{Diseño de la solución}

\subsection{Arquitectura general}

La solución estará compuesta por los siguientes módulos:

\begin{enumerate}
    \item Módulo de captura de imágenes.
    \item Módulo de procesamiento OCR.
    \item Motor de Inteligencia Artificial.
    \item API de integración.
    \item Plataforma web administrativa.
\end{enumerate}

\subsection{Tecnologías sugeridas}

\begin{itemize}
    \item OCR Tesseract.
    \item EasyOCR.
    \item PaddleOCR.
    \item FastAPI o Flask.
    \item PostgreSQL.
    \item Docker.
\end{itemize}

\section{Flujo operativo del sistema}

\begin{enumerate}
    \item El cajero toma una fotografía del voucher.
    \item La imagen es enviada al servidor IA.
    \item El OCR extrae el texto.
    \item El modelo IA identifica los campos importantes.
    \item El sistema valida formatos y duplicados.
    \item La información es registrada automáticamente.
    \item El voucher queda almacenado digitalmente.
\end{enumerate}

\section{Técnicas de Inteligencia Artificial empleadas}

\subsection{Visión por computador}

Utilizada para mejorar calidad de imágenes, detectar bordes, corregir inclinaciones y eliminar ruido.

\subsection{OCR Inteligente}

Permite convertir texto contenido en imágenes a texto digital editable.

\subsection{Machine Learning supervisado}

Se entrenará un modelo utilizando vouchers históricos para reconocer patrones de diferentes bancos e identificar errores comunes.

\section{Recursos tecnológicos necesarios}

\subsection{Alternativa 1: Implementación en VPS}

\begin{longtable}{|p{5cm}|p{8cm}|}
\hline
\textbf{Recurso} & \textbf{Recomendación} \\
\hline
CPU & 4 a 8 vCPU \\
\hline
RAM & 8 GB a 16 GB \\
\hline
Disco & 120 GB a 250 GB SSD \\
\hline
Sistema Operativo & Ubuntu Server 24.04 \\
\hline
GPU & NVIDIA T4 opcional \\
\hline
\end{longtable}

\subsection{Alternativa 2: Implementación local}

\begin{longtable}{|p{5cm}|p{8cm}|}
\hline
\textbf{Recurso} & \textbf{Recomendación} \\
\hline
Procesador & Intel Core i7 o Ryzen 7 \\
\hline
RAM & 16 GB \\
\hline
Disco & SSD 1 TB \\
\hline
GPU & NVIDIA RTX 3060 o superior \\
\hline
Sistema Operativo & Ubuntu o Windows 11 Pro \\
\hline
\end{longtable}

\section{Metodología de desarrollo}

\subsection{Metodología seleccionada: Scrum}

Se selecciona Scrum debido a la necesidad de iteraciones rápidas, retroalimentación continua y entrenamiento progresivo del modelo IA.

\subsection{Roles del proyecto}

\begin{longtable}{|p{5cm}|p{8cm}|}
\hline
\textbf{Rol} & \textbf{Función} \\
\hline
Product Owner & Define requerimientos \\
\hline
Scrum Master & Coordina metodología \\
\hline
Ingeniero IA & Entrena modelos \\
\hline
Backend Developer & APIs y lógica \\
\hline
Frontend Developer & Plataforma web \\
\hline
DevOps & Infraestructura y despliegue \\
\hline
QA Tester & Validaciones \\
\hline
Analista funcional & Procesos de negocio \\
\hline
\end{longtable}

\section{Seguridad y anonimización}

\subsection{Seguridad}

Se implementarán HTTPS, autenticación JWT, roles y permisos, logs de auditoría y firewalls.

\subsection{Anonimización}

Se ocultarán números completos de tarjeta, se cifrará información sensible y se aplicará retención limitada de imágenes.

\section{Indicadores KPI}

\begin{longtable}{|p{7cm}|p{5cm}|}
\hline
\textbf{KPI} & \textbf{Objetivo} \\
\hline
Precisión OCR & >95\% \\
\hline
Tiempo procesamiento & <5 segundos \\
\hline
Reducción digitación manual & 80\% \\
\hline
Errores de captura & <2\% \\
\hline
Disponibilidad sistema & 99\% \\
\hline
\end{longtable}

\section{Conclusiones}

La implementación de un sistema basado en Inteligencia Artificial para reconocimiento automático de vouchers representa una solución estratégica para modernizar los procesos de cuadre de caja en tiendas y puntos de venta.

La combinación de OCR, visión por computador y aprendizaje automático permite automatizar tareas repetitivas y mejorar significativamente la eficiencia operativa.

\end{document}
